% GENERATED FILE. Edit canonical lessons, then run npm run generate:books.
% canonical-source-hash: fnv1a64:a00681987cfeb7c6
% canonical-lessons: LA-C535-amplector, LA-C535-complector, LA-C535-amans, LA-C535-amator, LA-C535-procus

\chapter{Lovers}
\label{ch:la-lovers}
\hlchaptermodality{535}

\begin{chapteropening}
\textbf{By the end of this chapter:} \emph{I can talk about lovers, admirers and suitors.}

Complete the last lesson of chapter 535: I can talk about lovers, admirers and suitors.
\end{chapteropening}

\section[amplector, amplectī]{amplector, amplectī — to embrace, hug}

\label{lesson:LA-C535-amplector}



\begin{quote}
Before the new one: say the Latin for a kiss on the lips, then the Latin for an embrace.
\end{quote}

\subsection*{You'll want to know: amplector, amplectī}
\textbf{amplector, amplectī} — \textquotedblleft{}to embrace, hug\textquotedblright{}.

\begin{quote}
Pater fīlium redeuntem amplectitur.
\end{quote}

\begin{quote}
\emph{The father embraces his son as he comes home.}
\end{quote}

Deponent: \emph{amplector, amplectī, amplexus sum}: \emph{ambi-} 'around' + \emph{plectō} 'to twine'. Also to welcome an idea.

\begin{cousinweb}
\textbf{amplector} is built on \emph{ambi-} \textquotedblleft{}around\textquotedblright{} and \emph{plector}, from \emph{plectō} \textquotedblleft{}to twine, plait\textquotedblright{} (EtymDB): to wind the arms around.
\end{cousinweb}

\begin{grammarlens}[title={What you've built}]
One more word for telling, explaining and describing at B1.
\end{grammarlens}

\subsection*{Your turn}
\begin{itemize}
\raggedright
  \item \emph{Say it:} \emph{amplector, amplectī}
  \item \emph{Say it:} \emph{amplector, amplectī}, once more
  \item \emph{Say it:} a sentence of your own with \emph{amplector, amplectī}
  \item \emph{Recall:} say the Latin for a kiss on the lips, then the Latin for an embrace, then say \emph{amplector, amplectī} again
\end{itemize}

\begin{tcolorbox}[breakable,colback=teal!4,colframe=teal!35!black,title={Before you move on}]
What does amplector, amplectī mean? (\textquotedblleft{}To embrace, hug\textquotedblright{}.) Say it once more.
\end{tcolorbox}
\section[complector, complectī]{complector, complectī — to clasp, embrace; to include}

\label{lesson:LA-C535-complector}



\begin{quote}
Before the new one: say the Latin for an embrace, then the Latin for to embrace.
\end{quote}

\subsection*{You'll want to know: complector, complectī}
\textbf{complector, complectī} — \textquotedblleft{}to clasp, embrace; to include\textquotedblright{}.

\begin{quote}
Amīcī post longum tempus sē complectuntur.
\end{quote}

\begin{quote}
\emph{The friends embrace after a long time.}
\end{quote}

Deponent: \emph{complector, complectī, complexus sum}. Also to take in, to include. English \emph{complex} comes from its participle.

\begin{cousinweb}
\textbf{complector} is built on \emph{con-} and \emph{plectō} \textquotedblleft{}to twine\textquotedblright{} (EtymDB; Lewis \& Short). English \emph{complex} comes from its participle \emph{complexus}.
\end{cousinweb}

\begin{grammarlens}[title={What you've built}]
One more word for telling, explaining and describing at B1.
\end{grammarlens}

\subsection*{Your turn}
\begin{itemize}
\raggedright
  \item \emph{Say it:} \emph{complector, complectī}
  \item \emph{Say it:} \emph{complector, complectī}, once more
  \item \emph{Say it:} a sentence of your own with \emph{complector, complectī}
  \item \emph{Recall:} say the Latin for an embrace, then the Latin for to embrace, then say \emph{complector, complectī} again
\end{itemize}

\begin{tcolorbox}[breakable,colback=teal!4,colframe=teal!35!black,title={Before you move on}]
What does complector, complectī mean? (\textquotedblleft{}To clasp, embrace; to include\textquotedblright{}.) Say it once more.
\end{tcolorbox}
\section[amāns, amantis]{amāns, amantis — a lover}

\label{lesson:LA-C535-amans}



\begin{quote}
Before the new one: say the Latin for to embrace, then the Latin for to clasp.
\end{quote}

\subsection*{You'll want to know: amāns, amantis}
\textbf{amāns, amantis} — \textquotedblleft{}a lover\textquotedblright{}.

\begin{quote}
Amāns sub fenestrā puellae cantat.
\end{quote}

\begin{quote}
\emph{The lover sings under the girl's window.}
\end{quote}

The present participle of \emph{amō}, used as a noun. \emph{Amantēs} — lovers.

\begin{cousinweb}
\textbf{amāns} is the present participle of \emph{amō} \textquotedblleft{}to love\textquotedblright{}, used as a noun (Lewis \& Short). English \emph{amorous} belongs to the same family.
\end{cousinweb}

\begin{grammarlens}[title={What you've built}]
One more word for telling, explaining and describing at B1.
\end{grammarlens}

\subsection*{Your turn}
\begin{itemize}
\raggedright
  \item \emph{Say it:} \emph{amāns, amantis}
  \item \emph{Say it:} \emph{amāns, amantis}, once more
  \item \emph{Say it:} a sentence of your own with \emph{amāns, amantis}
  \item \emph{Recall:} say the Latin for to embrace, then the Latin for to clasp, then say \emph{amāns, amantis} again
\end{itemize}

\begin{tcolorbox}[breakable,colback=teal!4,colframe=teal!35!black,title={Before you move on}]
What does amāns, amantis mean? (\textquotedblleft{}A lover\textquotedblright{}.) Say it once more.
\end{tcolorbox}
\section[amātor, amātōris]{amātor, amātōris — an admirer, lover}

\label{lesson:LA-C535-amator}



\begin{quote}
Before the new one: say the Latin for to clasp, then the Latin for a lover.
\end{quote}

\subsection*{You'll want to know: amātor, amātōris}
\textbf{amātor, amātōris} — \textquotedblleft{}an admirer, lover\textquotedblright{}.

\begin{quote}
Amātor puellae flōrēs mittit.
\end{quote}

\begin{quote}
\emph{Her admirer sends the girl flowers.}
\end{quote}

Masculine, third declension. Also a lover of things: \emph{amātor librōrum}. English \emph{amateur} comes from it through French.

\begin{cousinweb}
\textbf{amātor} is built on \emph{amō} \textquotedblleft{}to love\textquotedblright{} with \emph{-tor} (Lewis \& Short; EtymDB). English \emph{amateur}, one who does a thing for love, comes from it through French.
\end{cousinweb}

\begin{grammarlens}[title={What you've built}]
One more word for telling, explaining and describing at B1.
\end{grammarlens}

\subsection*{Your turn}
\begin{itemize}
\raggedright
  \item \emph{Say it:} \emph{amātor, amātōris}
  \item \emph{Say it:} \emph{amātor, amātōris}, once more
  \item \emph{Say it:} a sentence of your own with \emph{amātor, amātōris}
  \item \emph{Recall:} say the Latin for to clasp, then the Latin for a lover, then say \emph{amātor, amātōris} again
\end{itemize}

\begin{tcolorbox}[breakable,colback=teal!4,colframe=teal!35!black,title={Before you move on}]
What does amātor, amātōris mean? (\textquotedblleft{}An admirer, lover\textquotedblright{}.) Say it once more.
\end{tcolorbox}
\section[procus, procī]{procus, procī — a suitor}

\label{lesson:LA-C535-procus}



\begin{quote}
Before the new one: say the Latin for a lover, then the Latin for an admirer.
\end{quote}

\subsection*{You'll want to know: procus, procī}
\textbf{procus, procī} — \textquotedblleft{}a suitor\textquotedblright{}.

\begin{quote}
Multī procī fīliam rēgis petunt.
\end{quote}

\begin{quote}
\emph{Many suitors seek the king's daughter.}
\end{quote}

Masculine, second declension; usually plural. A word of epic and comedy.

\begin{cousinweb}
\textbf{procus} is inherited from a Proto-Indo-European root meaning \textquotedblleft{}to ask\textquotedblright{} (EtymDB), the same root as \emph{precor} \textquotedblleft{}to pray\textquotedblright{}: a suitor is one who asks.
\end{cousinweb}

\begin{grammarlens}[title={What you've built}]
One more word for telling, explaining and describing at B1.
\end{grammarlens}

\subsection*{Your turn}
\begin{itemize}
\raggedright
  \item \emph{Say it:} \emph{procus, procī}
  \item \emph{Say it:} \emph{procus, procī}, once more
  \item \emph{Say it:} a sentence of your own with \emph{procus, procī}
  \item \emph{Recall:} say the Latin for a lover, then the Latin for an admirer, then say \emph{procus, procī} again
\end{itemize}

\begin{tcolorbox}[breakable,colback=teal!4,colframe=teal!35!black,title={Before you move on}]
What does procus, procī mean? (\textquotedblleft{}A suitor\textquotedblright{}.) Say it once more.
\end{tcolorbox}
